\documentclass[11pt, a4paper]{amsart}
\usepackage{amsmath, amssymb, amsthm, amsfonts}
\usepackage{mathtools}
\usepackage[margin=1in]{geometry}
\usepackage{hyperref}
\usepackage{tikz-cd}

\newtheorem{theorem}{Theorem}[section]
\newtheorem{lemma}[theorem]{Lemma}
\newtheorem{proposition}[theorem]{Proposition}
\newtheorem{corollary}[theorem]{Corollary}
\newtheorem{fact}[theorem]{Standard Fact}
\theoremstyle{definition}
\newtheorem{definition}[theorem]{Definition}
\newtheorem{remark}[theorem]{Remark}
\newtheorem{example}[theorem]{Example}
\newtheorem{notation}[theorem]{Notation}

\newcommand{\R}{\mathbb{R}}
\newcommand{\Q}{\mathbb{Q}}
\newcommand{\C}{\mathbb{C}}
\newcommand{\Z}{\mathbb{Z}}
\newcommand{\Fix}{\operatorname{Fix}}
\newcommand{\id}{\operatorname{id}}
\newcommand{\Hom}{\operatorname{Hom}}
\newcommand{\CP}{\mathbb{CP}}
\newcommand{\im}{\operatorname{im}}

\title{Cascade--Classical Correspondence I:\\
  Sigma-Fixed Rationality over $\mathbb{Q}(\varphi)$}
\author{}
\date{}

\begin{document}

\maketitle

\begin{abstract}
Let $K = \Q(\varphi)$ denote the quadratic extension of $\Q$ by the
golden ratio $\varphi = (1+\sqrt{5})/2$, and let
$\sigma \colon K \to K$ denote its unique non-trivial $\Q$-automorphism,
$\sigma(\varphi) = 1 - \varphi$. This paper proves two elementary but
load-bearing results about the coefficientwise $\sigma$-action on
scalar extensions $V_K := V \otimes_\Q K$:
\begin{enumerate}
\item (Fixed-part theorem.) For every finite-dimensional $\Q$-vector
  space $V$, the fixed subspace of $\sigma_V := \id_V \otimes \sigma$
  on $V_K$ is canonically $V \otimes 1 \cong V$
  (Theorem~\ref{thm:fixed-part}).
\item (Functoriality of $\Pi_\sigma$.) The $\Q$-linear projector
  $\Pi_\sigma := (\id + \sigma_V)/2$ commutes with the scalar
  extension of every $\Q$-linear map
  (Theorem~\ref{thm:pisigma-functorial}), and in particular with
  every cohomological pullback or pushforward defined over $\Q$
  (Corollaries~\ref{cor:cohomological} and
  \ref{cor:cohomological-pushforward}).
\end{enumerate}

These results are \emph{not} the Hodge Conjecture, nor do they
reduce it. They are the algebraic content of one specific
``$\sigma$-fixed equals rational'' step used in downstream
papers. The paper is deliberately narrow: no cascade-embeddability
claim is made, no universality statement is asserted, no geometric
or dynamical $\sigma$-action is constructed. Theorem inputs are
standard algebra only; no ``formal framework'' or exploratory
cascade note is used as theorem support. This is the first of
thirteen planned infrastructure papers and exists to provide a
stable citation target for the Foundations paper and the Hodge
companion manuscript.
\end{abstract}

\tableofcontents

%=====================================================================
\section{Introduction and Source Audit}
\label{sec:intro}
%=====================================================================

\subsection{Where $\sigma$ comes from in the wider programme}

In the wider cascade programme, $\sigma$ arises as a natural
involution on golden-field arithmetic: the non-trivial element of
$\operatorname{Gal}(\Q(\varphi)/\Q)$. This provenance (historically
from icosian arithmetic used in the pairing of the 600-cell with
its dual 120-cell) is not needed for any proof in the present
paper. The only fact we require is that $K := \Q(\varphi)$ is a
quadratic field extension of $\Q$ with a unique non-trivial
$\Q$-automorphism $\sigma$. Everything that follows is standard
algebra.

\subsection{Main results}

Theorem~\ref{thm:fixed-part} establishes that the $\sigma$-fixed
subspace of the scalar extension $V_K = V \otimes_\Q K$ is, under
the canonical inclusion $V \hookrightarrow V_K$, $v \mapsto
v \otimes 1$, equal to the image of $V$. Theorem~\ref{thm:pisigma-functorial}
establishes that the $\Q$-linear projector
$\Pi_\sigma := (\id + \sigma_V)/2 \colon V_K \to V_K$ commutes with
the scalar extension $T_K := T \otimes \id_K$ of every $\Q$-linear
map $T \colon V \to W$, i.e.\ the diagram
\[
\begin{tikzcd}
  V_K \arrow[r, "T_K"] \arrow[d, "\Pi_\sigma"'] & W_K \arrow[d, "\Pi_\sigma"] \\
  V_K \arrow[r, "T_K"'] & W_K
\end{tikzcd}
\]
commutes. Corollaries~\ref{cor:cohomological} and
\ref{cor:cohomological-pushforward} specialize this to
cohomological pullback and pushforward maps defined over $\Q$.

\subsection{What this paper does not prove}
\label{subsec:does-not-prove}

To forestall any reader mis-reading this as a Hodge paper, we
record the exclusions explicitly.

\begin{enumerate}
\item This paper does not prove the Hodge Conjecture or any part
  of it.
\item This paper does not prove that every Hodge class is
  $\sigma$-fixed.
\item This paper does not prove that every $\sigma$-fixed class is
  algebraic.
\item This paper does not claim any smooth projective variety is
  ``cascade-embeddable'' in any sense.
\item This paper does not propose a non-coefficientwise
  $\sigma$-action on varieties, cohomology classes, or cascade
  configurations.
\item This paper does not claim any relation between $\Pi_\sigma$
  and universality, motives, the Langlands programme, or any
  Millennium conjecture conclusion.
\end{enumerate}

What we do prove is a pair of clean algebraic lemmas about
coefficient extensions from $\Q$ to $\Q(\varphi)$. The downstream
papers (Foundations; the Hodge companion manuscript) will cite
these lemmas exactly as stated, and any step from
``$\sigma$-fixed'' to ``algebraic'' will be done in those papers,
not here.

\subsection{Source audit}

The only theorem inputs to this paper are elementary algebra and,
in \S\ref{sec:examples}, standard rational-cohomology conventions
for complex projective spaces and complex abelian varieties. No
internal cascade-derivation source is used as a theorem input. In
particular: the pregeometry sections P4--P8, the foundations
sections F2, F3, F6, F7, F8, the GR sections C2.bis and C3.bis,
the dynamics sections D5--D10, and the earlier Hodge draft are
\emph{not} used as proof support anywhere in this paper. Those
sources are either formal frameworks or exploratory notes in their
own source documents, and the purpose of this infrastructure paper
is precisely to \emph{replace} such dependence in downstream work.

\subsection{Organization}

Section~\ref{sec:prelim} develops algebraic preliminaries: the
field $K = \Q(\varphi)$, the involution $\sigma$, and the scalar
extension $V_K$. Section~\ref{sec:fixed-part} states and proves
the fixed-part theorem (Theorem~\ref{thm:fixed-part}).
Section~\ref{sec:pisigma} defines the projector $\Pi_\sigma$,
proves its basic properties, and establishes functoriality
(Theorem~\ref{thm:pisigma-functorial} and
Corollaries~\ref{cor:cohomological} and
\ref{cor:cohomological-pushforward}). Section~\ref{sec:examples}
works out three sanity-check examples.
Section~\ref{sec:handoff} states the citation contract for
downstream papers. Appendix~\ref{app:notation} is a concordance of
symbols and formulas.

%=====================================================================
\section{Algebraic Preliminaries: $\Q(\varphi)$, $\sigma$, and Scalar Extension}
\label{sec:prelim}
%=====================================================================

\subsection{The field $\Q(\varphi)$ and its Galois involution}

\begin{definition}[The field $K$]
\label{def:K}
Let
\[
  K := \Q(\varphi) = \Q[x]/(x^2 - x - 1),
\]
where $\varphi = (1 + \sqrt{5})/2$ is the golden ratio, a root of
$x^2 - x - 1 = 0$. As a $\Q$-vector space, $K$ is two-dimensional,
with basis $\{1, \varphi\}$.
\end{definition}

\begin{definition}[The involution $\sigma$]
\label{def:sigma}
The \emph{Galois involution} $\sigma \colon K \to K$ is the unique
non-trivial $\Q$-automorphism of $K$. Explicitly,
\[
  \sigma(a + b \varphi) := a + b \bar\varphi,
  \qquad a, b \in \Q,
\]
where $\bar\varphi := 1 - \varphi$ is the other root of
$x^2 - x - 1$. Equivalently, $\sigma(\varphi) = 1 - \varphi =
\bar\varphi$.
\end{definition}

\begin{remark}[Uniqueness of $\sigma$]
The minimal polynomial $x^2 - x - 1$ has exactly two roots in $K$,
namely $\varphi$ and $1 - \varphi$. Any $\Q$-automorphism of $K$
permutes these roots. The identity sends $\varphi \mapsto \varphi$;
the non-identity automorphism sends $\varphi \mapsto 1 - \varphi$.
This is the content of \cite[Ch.~V \S 1]{LangAlgebra} specialized to
the quadratic extension $K/\Q$.
\end{remark}

\begin{notation}[$\varphi$-basis vs.\ $\delta$-basis]
\label{not:bases}
We use two coordinate systems on $K$:
\begin{itemize}
\item The \emph{$\varphi$-basis}: $\{1, \varphi\}$. In this basis,
  $\sigma(a + b\varphi) = (a + b) - b\varphi$.
\item The \emph{$\delta$-basis}: $\{1, \delta\}$, where $\delta :=
  2\varphi - 1 = \sqrt{5}$. In this basis,
  $\sigma(a + b\delta) = a - b\delta$, the sign-flip involution.
\end{itemize}
The $\delta$-basis is simpler for the proof of the fixed-part
theorem, because $\sigma$ acts diagonally as $+1$ on the first
basis vector and $-1$ on the second.
\end{notation}

\begin{lemma}[Coordinate change]
\label{lem:coord-change}
The two bases are related by $\delta = 2\varphi - 1$ and $\varphi
= (1 + \delta)/2$. In particular, a general element of $K$ can be
written either as $a + b\varphi$ or as $A + B\delta$, with $A =
a + b/2$ and $B = b/2$. The involution $\sigma$ acts on the
$\delta$-basis by
\[
  \sigma(A + B\delta) = A - B\delta,
  \qquad A, B \in \Q.
\]
\end{lemma}

\begin{proof}
$\sigma(\delta) = \sigma(2\varphi - 1) = 2\sigma(\varphi) - 1 =
2(1-\varphi) - 1 = 1 - 2\varphi = -(2\varphi - 1) = -\delta$.
Therefore $\sigma(A + B\delta) = A + B\sigma(\delta) = A - B\delta$
by $\Q$-linearity of $\sigma$.
\end{proof}

\begin{lemma}[Direct-sum decomposition of $K$]
\label{lem:K-decomp}
As a $\Q$-vector space,
\[
  K = \Q \oplus \Q\delta,
\]
and the involution $\sigma \colon K \to K$ acts as $(+1, -1)$ with
respect to this decomposition. The $\sigma$-fixed subspace of $K$
is $\Fix(\sigma) = \Q$, and the $\sigma$-anti-fixed subspace is
$\Q\delta$.
\end{lemma}

\begin{proof}
By Definition~\ref{def:K}, $K$ has $\Q$-dimension $2$. The vectors
$1$ and $\delta = 2\varphi - 1$ are $\Q$-linearly independent
(since $\delta \notin \Q$; equivalently, $\sqrt{5}$ is irrational).
Hence $K = \Q \cdot 1 \oplus \Q \cdot \delta$ as a $\Q$-vector space.
The $\sigma$-action was computed in Lemma~\ref{lem:coord-change}:
$\sigma$ fixes $1$ and sends $\delta \mapsto -\delta$.
Hence the $\sigma$-eigenspaces are $\Q$ (eigenvalue $+1$) and
$\Q\delta$ (eigenvalue $-1$).
\end{proof}

\subsection{Scalar extension $V_K$ and the coefficientwise $\sigma$-action}

\begin{definition}[Scalar extension]
\label{def:VK}
Let $V$ be a finite-dimensional $\Q$-vector space. The
\emph{scalar extension of $V$ to $K$} is the $K$-vector space
\[
  V_K := V \otimes_\Q K,
\]
with $K$-action $\lambda \cdot (v \otimes \mu) := v \otimes
(\lambda \mu)$ for $\lambda, \mu \in K$ and $v \in V$. A
$\Q$-basis $\{e_1, \ldots, e_n\}$ of $V$ gives a $K$-basis
$\{e_i \otimes 1\}_{i=1}^n$ of $V_K$.
\end{definition}

\begin{definition}[Inclusion $V \hookrightarrow V_K$]
\label{def:inclusion}
The $\Q$-linear map
\[
  \iota_V \colon V \hookrightarrow V_K,
  \qquad v \mapsto v \otimes 1,
\]
is injective (since $1 \neq 0$ in $K$) and $\Q$-linear, but not
$K$-linear (its image is not a $K$-subspace of $V_K$). We often
identify $v \in V$ with its image $v \otimes 1 \in V_K$ by this
inclusion.
\end{definition}

\begin{definition}[Coefficientwise $\sigma$-action]
\label{def:sigmaV}
The \emph{coefficientwise $\sigma$-action} on $V_K$ is
\[
  \sigma_V := \id_V \otimes \sigma \colon V_K \to V_K,
  \qquad (\id_V \otimes \sigma)(v \otimes \lambda)
  = v \otimes \sigma(\lambda),
\]
extended $\Q$-bilinearly. Equivalently, if $\omega = \sum_i v_i
\otimes \lambda_i \in V_K$ with $v_i \in V$ and $\lambda_i \in K$,
then $\sigma_V(\omega) = \sum_i v_i \otimes \sigma(\lambda_i)$.
\end{definition}

\begin{remark}[$\Q$-linear, $K$-semilinear]
\label{rmk:semilinear}
$\sigma_V$ is $\Q$-linear by construction (both $\id_V$ and $\sigma$
are). However, $\sigma_V$ is \emph{not} $K$-linear: for $\lambda
\in K$ and $\omega \in V_K$,
\[
  \sigma_V(\lambda \cdot \omega) = \sigma(\lambda) \cdot \sigma_V(\omega),
\]
which equals $\lambda \sigma_V(\omega)$ only when $\sigma(\lambda)
= \lambda$, i.e.\ $\lambda \in \Q$. In the language of
\cite[Ch.~XVI \S 7]{LangAlgebra}, $\sigma_V$ is \emph{$\sigma$-semilinear}:
it is additive, and $\sigma_V(\lambda \omega) = \sigma(\lambda)
\sigma_V(\omega)$.
\end{remark}

\begin{proposition}[$\delta$-basis decomposition of $V_K$]
\label{prop:VK-decomp}
Let $V$ be a finite-dimensional $\Q$-vector space. Then as a
$\Q$-vector space,
\[
  V_K = V \oplus V\delta,
\]
where $V\delta := V \otimes_\Q (\Q\delta) = \{v \otimes \delta : v
\in V\}$. With respect to this decomposition, every element
$\omega \in V_K$ has a unique expression $\omega = v_0 + v_1\delta$
with $v_0, v_1 \in V$, and
\[
  \sigma_V(v_0 + v_1\delta) = v_0 - v_1\delta.
\]
\end{proposition}

\begin{proof}
By Lemma~\ref{lem:K-decomp}, $K = \Q \oplus \Q\delta$ as a
$\Q$-vector space. Tensoring over $\Q$ with $V$ (which preserves
direct sums) gives
\[
  V \otimes_\Q K = (V \otimes_\Q \Q) \oplus (V \otimes_\Q \Q\delta)
  = V \oplus V\delta,
\]
where the first summand identifies with $V$ via $v \otimes 1
\leftrightarrow v$. Uniqueness of the decomposition follows from
the direct-sum property. The $\sigma_V$-action on each summand is:
$\sigma_V(v \otimes 1) = v \otimes \sigma(1) = v \otimes 1$ and
$\sigma_V(v \otimes \delta) = v \otimes \sigma(\delta) = v \otimes
(-\delta) = -v \otimes \delta$. Linear extension gives the stated
formula.
\end{proof}

%=====================================================================
\section{The Fixed-Part Theorem}
\label{sec:fixed-part}
%=====================================================================

\begin{definition}[Fixed subspace]
\label{def:fix}
For $\sigma_V \colon V_K \to V_K$ as in Definition~\ref{def:sigmaV},
the \emph{$\sigma_V$-fixed subspace} is
\[
  \Fix(\sigma_V; V_K) := \{\omega \in V_K : \sigma_V(\omega) = \omega\}.
\]
This is a $\Q$-linear subspace of $V_K$ (intersection of the
kernel of the $\Q$-linear map $\sigma_V - \id$ with $V_K$).
\end{definition}

\begin{theorem}[Fixed-part theorem]
\label{thm:fixed-part}
Let $V$ be a finite-dimensional $\Q$-vector space. Then
\[
  \Fix(\sigma_V; V_K) = V \otimes 1 = \iota_V(V),
\]
which is canonically identified with $V$ under the inclusion
$\iota_V \colon V \hookrightarrow V_K$ of
Definition~\ref{def:inclusion}.
\end{theorem}

\begin{proof}
By Proposition~\ref{prop:VK-decomp}, every $\omega \in V_K$ has a
unique expression $\omega = v_0 + v_1 \delta$ with $v_0, v_1 \in V$,
and $\sigma_V(\omega) = v_0 - v_1 \delta$. The equation
$\sigma_V(\omega) = \omega$ is equivalent to
\[
  v_0 - v_1 \delta = v_0 + v_1 \delta,
\]
i.e.\ $-v_1 \delta = v_1 \delta$, i.e.\ $2 v_1 \delta = 0$ inside
$V_K$.

We must now deduce $v_1 = 0$. By
Proposition~\ref{prop:VK-decomp}, $V_K = V \oplus V\delta$ as an
internal direct sum of $\Q$-vector spaces, so the decomposition
$\omega = v_0 + v_1\delta$ with $v_0, v_1 \in V$ is unique. Since
$\operatorname{char} \Q = 0$, multiplication by $2$ is a $\Q$-linear
automorphism of $V_K$, so $2 v_1 \delta = 0$ in $V_K$ is equivalent
to $v_1 \delta = 0$ in $V_K$. Uniqueness of the $V \oplus V\delta$
decomposition applied to the zero element then forces $v_1 = 0$ in
$V$.

Hence
\[
  \omega \in \Fix(\sigma_V; V_K) \iff v_1 = 0 \iff \omega = v_0 \in V \otimes 1.
\]

Conversely, every $\omega = v_0 \in V \otimes 1$ satisfies
$\sigma_V(\omega) = v_0 - 0 \cdot \delta = v_0 = \omega$.

Therefore $\Fix(\sigma_V; V_K) = V \otimes 1$. The identification
with $V$ is via the inclusion $\iota_V$, which is injective and
$\Q$-linear (Definition~\ref{def:inclusion}). This gives a
canonical $\Q$-linear isomorphism $\Fix(\sigma_V; V_K) \cong V$.
\end{proof}

\begin{remark}[Why this is not Hodge]
\label{rmk:not-hodge}
Theorem~\ref{thm:fixed-part} says that the coefficientwise
$\sigma$-fixed part of a $\Q(\varphi)$-extended vector space
recovers the original $\Q$-vector space. This is a purely
algebraic fact about scalar extension in a quadratic Galois
extension of characteristic-zero fields. It has no geometric,
topological, or algebraic-geometric content on its own. In
particular, it says nothing about Hodge classes, algebraic cycles,
or any other Hodge-theoretic notion.
\end{remark}

\begin{corollary}[$\sigma_V$-eigenspace decomposition]
\label{cor:eigenspaces}
Over $\Q$,
\[
  V_K = \Fix(\sigma_V; V_K) \oplus \operatorname{AntiFix}(\sigma_V; V_K),
\]
where $\operatorname{AntiFix}(\sigma_V; V_K) := \{\omega \in V_K :
\sigma_V(\omega) = -\omega\} = V\delta$. Both summands have
$\Q$-dimension $\dim_\Q V$.
\end{corollary}

\begin{proof}
Proposition~\ref{prop:VK-decomp} gives $V_K = V \oplus V\delta$,
and the same argument as in the proof of
Theorem~\ref{thm:fixed-part} gives $\operatorname{AntiFix} = V\delta$.
Each summand has $\Q$-dimension equal to $\dim_\Q V$.
\end{proof}

%=====================================================================
\section{The Projector $\Pi_\sigma$ and Its Functoriality}
\label{sec:pisigma}
%=====================================================================

\subsection{Definition and basic properties}

\begin{definition}[The projector $\Pi_\sigma$]
\label{def:pisigma}
Let $V$ be a finite-dimensional $\Q$-vector space. The
\emph{$\sigma$-fixed-part projector} on $V_K$ is
\[
  \Pi_\sigma \colon V_K \to V_K,
  \qquad \Pi_\sigma(\omega) := \frac{\omega + \sigma_V(\omega)}{2}.
\]
\end{definition}

\begin{proposition}[Basic properties of $\Pi_\sigma$]
\label{prop:pisigma-basic}
For every finite-dimensional $\Q$-vector space $V$, the map
$\Pi_\sigma \colon V_K \to V_K$ satisfies:
\begin{enumerate}
\item $\Pi_\sigma$ is $\Q$-linear.
\item $\Pi_\sigma$ is generally not $K$-linear (see
  Remark~\ref{rmk:pisigma-not-Klinear}).
\item $\Pi_\sigma \circ \Pi_\sigma = \Pi_\sigma$ (idempotence).
\item $\Pi_\sigma \circ \sigma_V = \Pi_\sigma$.
\item $\im(\Pi_\sigma) = \Fix(\sigma_V; V_K)$.
\item $\ker(\Pi_\sigma) = \operatorname{AntiFix}(\sigma_V; V_K) = V\delta$.
\item In $\delta$-basis coordinates, $\Pi_\sigma(v_0 + v_1\delta)
  = v_0$.
\end{enumerate}
\end{proposition}

\begin{proof}
\emph{(1)} $\Pi_\sigma$ is $\Q$-linear because $\sigma_V$ is
$\Q$-linear (Remark~\ref{rmk:semilinear}), $\id$ is $\Q$-linear, and
averaging is a $\Q$-linear operation.

\emph{(2)} This is postponed to Remark~\ref{rmk:pisigma-not-Klinear}
after we have stated it.

\emph{(3)} Applying $\sigma_V$ to an element of
$\Fix(\sigma_V; V_K)$ returns the same element, so
\begin{align*}
  \sigma_V(\Pi_\sigma(\omega))
  &= \sigma_V\!\left(\frac{\omega + \sigma_V(\omega)}{2}\right)
  = \frac{\sigma_V(\omega) + \sigma_V^2(\omega)}{2}
  = \frac{\sigma_V(\omega) + \omega}{2}
  = \Pi_\sigma(\omega).
\end{align*}
Here $\sigma_V^2 = \id$ since $\sigma^2 = \id$ on $K$ and
$\sigma_V = \id_V \otimes \sigma$. So $\Pi_\sigma(\omega) \in
\Fix(\sigma_V; V_K)$, and applying $\Pi_\sigma$ again gives
\begin{align*}
  \Pi_\sigma(\Pi_\sigma(\omega))
  &= \frac{\Pi_\sigma(\omega) + \sigma_V(\Pi_\sigma(\omega))}{2}
  = \frac{\Pi_\sigma(\omega) + \Pi_\sigma(\omega)}{2}
  = \Pi_\sigma(\omega).
\end{align*}

\emph{(4)} Using $\sigma_V^2 = \id$,
\[
  \Pi_\sigma(\sigma_V(\omega))
  = \frac{\sigma_V(\omega) + \sigma_V^2(\omega)}{2}
  = \frac{\sigma_V(\omega) + \omega}{2}
  = \Pi_\sigma(\omega).
\]

\emph{(5)} By \emph{(3)}, $\im(\Pi_\sigma) \subseteq
\Fix(\sigma_V; V_K)$. Conversely, if $\omega \in \Fix(\sigma_V; V_K)$,
then $\Pi_\sigma(\omega) = (\omega + \omega)/2 = \omega$, so
$\omega \in \im(\Pi_\sigma)$. Hence $\im(\Pi_\sigma) = \Fix(\sigma_V;
V_K)$.

\emph{(6)} Using Proposition~\ref{prop:VK-decomp} with $\omega =
v_0 + v_1 \delta$:
\[
  \Pi_\sigma(\omega)
  = \frac{(v_0 + v_1 \delta) + (v_0 - v_1 \delta)}{2}
  = \frac{2 v_0}{2} = v_0.
\]
Hence $\Pi_\sigma(\omega) = 0$ iff $v_0 = 0$, i.e.\ $\omega \in
V\delta = \operatorname{AntiFix}(\sigma_V; V_K)$.

\emph{(7)} Immediate from the computation of $\Pi_\sigma(v_0 + v_1
\delta) = v_0$ above.

\emph{(2), postponed.} See Remark~\ref{rmk:pisigma-not-Klinear}.
\end{proof}

\begin{remark}[$\Pi_\sigma$ is not $K$-linear]
\label{rmk:pisigma-not-Klinear}
For $\lambda \in K \setminus \Q$ and $\omega \in V_K$,
\[
  \Pi_\sigma(\lambda \omega)
  = \frac{\lambda\omega + \sigma_V(\lambda\omega)}{2}
  = \frac{\lambda\omega + \sigma(\lambda) \sigma_V(\omega)}{2},
\]
while
\[
  \lambda \Pi_\sigma(\omega)
  = \frac{\lambda\omega + \lambda \sigma_V(\omega)}{2}.
\]
These are equal only if $\sigma(\lambda) = \lambda$ (i.e.\ $\lambda
\in \Q$) or $\sigma_V(\omega) = 0$. Concretely, for $\lambda =
\varphi$ and $\omega = v \otimes 1$,
\[
  \Pi_\sigma(\varphi (v \otimes 1))
  = \Pi_\sigma(v \otimes \varphi)
  = \frac{v \otimes \varphi + v \otimes \bar\varphi}{2}
  = v \otimes \tfrac{\varphi + \bar\varphi}{2}
  = v \otimes \tfrac{1}{2},
\]
since $\varphi + \bar\varphi = \varphi + (1-\varphi) = 1$, while
$\varphi \Pi_\sigma(v \otimes 1) = \varphi (v \otimes 1) = v \otimes
\varphi$. These differ: $v \otimes \tfrac{1}{2} \neq v \otimes
\varphi$ unless $v = 0$.

Hence $\Pi_\sigma$ is $\Q$-linear but not $K$-linear. This
distinction matters: any downstream usage of $\Pi_\sigma$ must
treat it as a $\Q$-linear projector on the underlying
$\Q$-vector-space structure of $V_K$, not as a $K$-linear operator.
\end{remark}

\subsection{Functoriality: the main theorem}

\begin{definition}[Scalar extension of a $\Q$-linear map]
\label{def:TK}
Let $T \colon V \to W$ be a $\Q$-linear map of finite-dimensional
$\Q$-vector spaces. The \emph{scalar extension of $T$ to $K$} is
\[
  T_K := T \otimes \id_K \colon V_K \to W_K,
  \qquad (T \otimes \id_K)(v \otimes \lambda) = T(v) \otimes \lambda.
\]
$T_K$ is $K$-linear (not just $\Q$-linear): $T_K(\lambda(v \otimes
\mu)) = T_K(v \otimes (\lambda\mu)) = T(v) \otimes (\lambda\mu) =
\lambda(T(v) \otimes \mu) = \lambda T_K(v \otimes \mu)$.
\end{definition}

\begin{theorem}[Functoriality of $\sigma_V$ and $\Pi_\sigma$]
\label{thm:pisigma-functorial}
Let $T \colon V \to W$ be a $\Q$-linear map of finite-dimensional
$\Q$-vector spaces. Then both $\sigma$ and $\Pi_\sigma$ are
functorial with respect to $T$: the squares
\[
\begin{tikzcd}
  V_K \arrow[r, "T_K"] \arrow[d, "\sigma_V"'] & W_K \arrow[d, "\sigma_W"] \\
  V_K \arrow[r, "T_K"'] & W_K
\end{tikzcd}
\qquad
\begin{tikzcd}
  V_K \arrow[r, "T_K"] \arrow[d, "\Pi_\sigma^V"'] & W_K \arrow[d, "\Pi_\sigma^W"] \\
  V_K \arrow[r, "T_K"'] & W_K
\end{tikzcd}
\]
commute, i.e.\
\[
  \sigma_W \circ T_K = T_K \circ \sigma_V
  \quad \text{and} \quad
  \Pi_\sigma^W \circ T_K = T_K \circ \Pi_\sigma^V.
\]
\end{theorem}

\begin{proof}
We prove commutation of the first square; the second follows by
averaging.

For any simple tensor $v \otimes \lambda \in V_K$:
\begin{align*}
  (\sigma_W \circ T_K)(v \otimes \lambda)
  &= \sigma_W(T(v) \otimes \lambda) && \text{by Def.~\ref{def:TK}} \\
  &= T(v) \otimes \sigma(\lambda) && \text{by Def.~\ref{def:sigmaV} (applied to $W$)} \\
  &= T_K(v \otimes \sigma(\lambda)) && \text{by Def.~\ref{def:TK}} \\
  &= T_K(\sigma_V(v \otimes \lambda)) && \text{by Def.~\ref{def:sigmaV} (applied to $V$)} \\
  &= (T_K \circ \sigma_V)(v \otimes \lambda).
\end{align*}

The two maps $\sigma_W \circ T_K$ and $T_K \circ \sigma_V$ agree on
simple tensors, and both are $\Q$-linear (as compositions of
$\Q$-linear maps). Since $V_K$ is generated as a $\Q$-vector space
by simple tensors $v \otimes \lambda$, the two maps agree on all of
$V_K$. This establishes $\sigma_W \circ T_K = T_K \circ \sigma_V$.

For the $\Pi_\sigma$-square:
\begin{align*}
  \Pi_\sigma^W \circ T_K
  &= \tfrac{1}{2}(\id_{W_K} + \sigma_W) \circ T_K \\
  &= \tfrac{1}{2}(T_K + \sigma_W \circ T_K) \\
  &= \tfrac{1}{2}(T_K + T_K \circ \sigma_V) && \text{(just proved)} \\
  &= T_K \circ \tfrac{1}{2}(\id_{V_K} + \sigma_V) \\
  &= T_K \circ \Pi_\sigma^V.
\end{align*}
Here we used $\Q$-linearity of all maps (and in particular the
fact that scalar multiplication by $1/2$ is $\Q$-linear).
\end{proof}

\begin{definition}[Cohomological map defined over $\Q$]
\label{def:defined-over-Q}
Throughout this definition and Corollaries~\ref{cor:cohomological}
and~\ref{cor:cohomological-pushforward}, $H^*(\cdot, \Q)$ denotes
singular cohomology with $\Q$-coefficients \cite[Ch.~3]{Hatcher};
the canonical identification $H^*(\cdot, K) \cong H^*(\cdot, \Q)
\otimes_\Q K$ is then the universal coefficient isomorphism
combined with flatness of $K$ over $\Q$, as cited. Applications to
other rational cohomology theories (de~Rham, \'etale, crystalline,
etc.) require the analogous rationality / scalar-extension
identification as an additional hypothesis specific to that theory,
supplied externally; the present definition does not claim those
identifications automatically.

For topological spaces $X, Y$, a cohomological map with
$K$-coefficients
\[
  g_K \colon H^*(Y, K) \longrightarrow H^*(X, K)
\]
is said to be \emph{defined over $\Q$} if there exists a
$\Q$-linear map $g \colon H^*(Y, \Q) \to H^*(X, \Q)$ such that
$g_K = g \otimes \id_K$ in the sense of
Definition~\ref{def:TK}, under the canonical identification
$H^*(\cdot, K) \cong H^*(\cdot, \Q) \otimes_\Q K$ of singular
cohomology.
\end{definition}

\begin{corollary}[Cohomological version: pullback]
\label{cor:cohomological}
Let $f \colon X \to Y$ be a continuous map of topological spaces.
The induced singular-cohomology pullback
$f^* \colon H^*(Y, \Q) \to H^*(X, \Q)$ is $\Q$-linear, so its
scalar extension $f^*_K := f^* \otimes \id_K \colon H^*(Y, K) \to
H^*(X, K)$ is a cohomological map defined over $\Q$ in the sense
of Definition~\ref{def:defined-over-Q}. (An analogous statement
for de~Rham, \'etale, crystalline, or other rational cohomology
theories requires the analogous scalar-extension identification
$H^*(\cdot, K) \cong H^*(\cdot, \Q) \otimes_\Q K$ for that theory
as an additional external hypothesis, supplied outside this paper;
cf.\ the final paragraph of Definition~\ref{def:defined-over-Q}.)
The diagrams
\[
\begin{tikzcd}
  H^*(Y, K) \arrow[r, "f^*_K"] \arrow[d, "\sigma_{H^*(Y, \Q)}"'] & H^*(X, K) \arrow[d, "\sigma_{H^*(X, \Q)}"] \\
  H^*(Y, K) \arrow[r, "f^*_K"'] & H^*(X, K)
\end{tikzcd}
\qquad
\begin{tikzcd}
  H^*(Y, K) \arrow[r, "f^*_K"] \arrow[d, "\Pi_\sigma"'] & H^*(X, K) \arrow[d, "\Pi_\sigma"] \\
  H^*(Y, K) \arrow[r, "f^*_K"'] & H^*(X, K)
\end{tikzcd}
\]
commute, i.e.\
$\sigma_{H^*(X, \Q)} \circ f^*_K = f^*_K \circ \sigma_{H^*(Y, \Q)}$
and $\Pi_\sigma \circ f^*_K = f^*_K \circ \Pi_\sigma$.
\end{corollary}

\begin{proof}
$f^*$ is $\Q$-linear by the definition of cohomology with
$\Q$-coefficients \cite[Ch.~3]{Hatcher}. Its scalar extension
$f^*_K = f^* \otimes \id_K$ is exactly the map $T_K$ of
Definition~\ref{def:TK} applied to $T = f^*$, so
$f^*_K$ is defined over $\Q$ in the sense of
Definition~\ref{def:defined-over-Q}. Applying
Theorem~\ref{thm:pisigma-functorial} with $V = H^*(Y, \Q)$, $W =
H^*(X, \Q)$, and $T = f^*$ gives commutation of both diagrams.
\end{proof}

\begin{corollary}[Cohomological version: pushforward]
\label{cor:cohomological-pushforward}
\emph{Hypothesis} (cohomological $\Q$-linear pushforward): suppose
a $\Q$-linear map
$f_* \colon H^*(X, \Q) \to H^*(Y, \Q)$ is given as an externally
supplied input. We do not assume or assert any particular
construction of $f_*$; the corollary is purely algebraic in the
given $\Q$-linear map. A standard example where such an $f_*$ is
supplied is when $X, Y$ are oriented closed manifolds and $f$ is a
smooth map, with $f_*$ defined via Poincar\'e duality and the
pullback of $f$ (see \cite[Ch.~3]{Hatcher} for the Poincar\'e-duality
construction of a singular-cohomology pushforward on oriented closed
manifolds); this example is motivational only and is not used as a
hypothesis of the corollary.

Under the above hypothesis, the scalar extension
$f_{*,K} := f_* \otimes \id_K \colon H^*(X, K) \to
H^*(Y, K)$ is a cohomological map defined over $\Q$
(Definition~\ref{def:defined-over-Q}), and the diagrams
\[
\begin{tikzcd}
  H^*(X, K) \arrow[r, "f_{*,K}"] \arrow[d, "\sigma_{H^*(X, \Q)}"'] & H^*(Y, K) \arrow[d, "\sigma_{H^*(Y, \Q)}"] \\
  H^*(X, K) \arrow[r, "f_{*,K}"'] & H^*(Y, K)
\end{tikzcd}
\qquad
\begin{tikzcd}
  H^*(X, K) \arrow[r, "f_{*,K}"] \arrow[d, "\Pi_\sigma"'] & H^*(Y, K) \arrow[d, "\Pi_\sigma"] \\
  H^*(X, K) \arrow[r, "f_{*,K}"'] & H^*(Y, K)
\end{tikzcd}
\]
commute. The corollary applies to any rational-cohomology
pushforward whose definition is $\Q$-linear; the proof does not
depend on the specific construction of $f_*$ beyond its
$\Q$-linearity.
\end{corollary}

\begin{proof}
By hypothesis $f_*$ is $\Q$-linear. Apply
Theorem~\ref{thm:pisigma-functorial} with $V = H^*(X, \Q)$, $W =
H^*(Y, \Q)$, and $T = f_*$.
\end{proof}

%=====================================================================
\section{Applications and Examples}
\label{sec:examples}
%=====================================================================

This section illustrates Theorem~\ref{thm:fixed-part},
Proposition~\ref{prop:pisigma-basic}, and
Theorem~\ref{thm:pisigma-functorial} on three concrete examples.
None of these examples is used as a theorem input elsewhere; they
are sanity checks that the projector $\Pi_\sigma$ acts as claimed.

\subsection{Complex projective space}

\begin{example}[$V = H^*(\CP^n, \Q)$]
\label{ex:cpn}
Let $n \geq 1$. By standard topology
\cite[Theorem~3.19]{Hatcher},
\[
  H^*(\CP^n, \Q) \cong \Q[h]/(h^{n+1}),
\]
where $h \in H^2(\CP^n, \Q)$ is the hyperplane class.
This is a $\Q$-vector space of dimension $n + 1$, with basis
$\{1, h, h^2, \ldots, h^n\}$.

Scalar extension to $K = \Q(\varphi)$ gives
\[
  H^*(\CP^n, K) \cong H^*(\CP^n, \Q) \otimes_\Q K
  \cong K[h]/(h^{n+1}),
\]
a $K$-vector space of $K$-dimension $n + 1$, or a $\Q$-vector space
of $\Q$-dimension $2(n+1)$.

By Theorem~\ref{thm:fixed-part} applied with $V = H^*(\CP^n, \Q)$,
\[
  \Fix\big(\sigma_{H^*(\CP^n, \Q)};\, H^*(\CP^n, K)\big)
  = H^*(\CP^n, \Q) \otimes 1 \cong H^*(\CP^n, \Q),
\]
a $\Q$-vector space of dimension $n + 1$. In coordinates, for
$\omega = \sum_{k=0}^n (a_k + b_k \delta) h^k$ with $a_k, b_k \in
\Q$, we have
$\Pi_\sigma(\omega) = \sum_{k=0}^n a_k h^k$. This recovers the
$\Q$-rational part of $\omega$. From this point on we write
$\Pi_\sigma$ and $\sigma$ without the underlying-space subscript
when the ambient $\Q$-vector space is clear from context; the
subscripted forms $\Pi_\sigma^V$ and $\sigma_V$ are always the
precise objects.
\end{example}

\subsection{Complex abelian variety}

\begin{example}[$V = H^*(A, \Q)$ for an abelian variety $A$]
\label{ex:abelian}
Let $A$ be a complex abelian variety of complex dimension $g$
(so topologically, $A$ is the quotient of $\C^g$ by a rank-$2g$
lattice; real dimension $2g$). By standard results
\cite[Ch.~2]{GriffithsHarris}, $H^1(A, \Q)$ is a $\Q$-vector space
of dimension $2g$, and
\[
  H^*(A, \Q) \cong \bigwedge^* H^1(A, \Q),
\]
with $H^k(A, \Q) \cong \bigwedge^k H^1(A, \Q)$ of $\Q$-dimension
$\binom{2g}{k}$.

Tensoring with $K$,
\[
  H^*(A, K) \cong \bigwedge_K^* \big(H^1(A, K)\big),
\]
where $H^1(A, K) = H^1(A, \Q) \otimes_\Q K$ has $K$-dimension $2g$.
Theorem~\ref{thm:fixed-part} applied with
$V = H^*(A, \Q)$ gives
\[
  \Fix\big(\sigma_{H^*(A, \Q)};\, H^*(A, K)\big)
  = H^*(A, \Q) \otimes 1 \cong H^*(A, \Q).
\]
Here $\Pi_\sigma$ acts on the $K$-coefficients of a chosen
$\Q$-basis expansion of classes in $H^*(A, K)$; no multiplicative
claim is made. In particular, $\Pi_\sigma$ is \emph{not} an algebra
homomorphism with respect to the cup product or the wedge-algebra
structure on $H^*(A, K)$: for instance, with $\omega_1 = \delta
\alpha$ and $\omega_2 = \delta \beta$ (where $\alpha, \beta \in
H^*(A, \Q) \otimes 1$), one has $\omega_1 \wedge \omega_2 =
\delta^2 (\alpha \wedge \beta) = 5 (\alpha \wedge \beta)$ (using
$\delta^2 = 5$), so $\Pi_\sigma(\omega_1 \wedge \omega_2) = 5
(\alpha \wedge \beta)$, whereas $\Pi_\sigma(\omega_1) \wedge
\Pi_\sigma(\omega_2) = 0 \wedge 0 = 0$.
\end{example}

\subsection{A map example: $\CP^m \hookrightarrow \CP^n$}

\begin{example}[Functoriality under linear embedding]
\label{ex:embedding}
Let $m < n$, and let $f \colon \CP^m \hookrightarrow \CP^n$ be a
linear embedding (e.g.\ $[x_0 : \cdots : x_m] \mapsto [x_0 :
\cdots : x_m : 0 : \cdots : 0]$). By standard topology, the
pullback $f^* \colon H^*(\CP^n, \Q) \to H^*(\CP^m, \Q)$ is the
$\Q$-algebra map sending the hyperplane class $h_n$ to the
hyperplane class $h_m$, and in general $h_n^k \mapsto h_m^k$ for
$0 \leq k \leq m$, with $h_n^k \mapsto 0$ for $m < k \leq n$ (since
$h_m^{m+1} = 0$ in $H^*(\CP^m, \Q)$).

By Corollary~\ref{cor:cohomological}, $f^*$ extends to
$f^*_K \colon K[h_n]/(h_n^{n+1}) \to K[h_m]/(h_m^{m+1})$, and
$\Pi_\sigma \circ f^*_K = f^*_K \circ \Pi_\sigma$.

We verify this directly on a basis element. Take $\omega =
h_n + \delta h_n^2 \in H^*(\CP^n, K)$, valid for $n \geq 2$. Then:
\[
  \Pi_\sigma(\omega) = h_n, \qquad
  f^*_K(\omega) = h_m + \delta h_m^2, \qquad
  f^*_K(\Pi_\sigma(\omega)) = h_m,
\]
and
\[
  \Pi_\sigma(f^*_K(\omega)) = \Pi_\sigma(h_m + \delta h_m^2)
  = h_m.
\]
So $\Pi_\sigma(f^*_K(\omega)) = f^*_K(\Pi_\sigma(\omega)) = h_m$,
confirming commutation.
\end{example}

\begin{remark}[What these examples do and do not show]
Examples~\ref{ex:cpn}, \ref{ex:abelian}, \ref{ex:embedding} are
sanity checks at the level of rational-cohomology bookkeeping.
They do not address whether any $\sigma$-fixed cohomology class is
algebraic, or whether $\Pi_\sigma$ plays any role in detecting
Hodge classes, or whether any variety is ``cascade-embeddable'' in
any sense. Those questions lie outside this paper by design.
\end{remark}

%=====================================================================
\section{Companion-Paper Citation Contract}
\label{sec:handoff}
%=====================================================================

This section states exactly what downstream papers may cite from
this paper and what they must independently prove or assume.

\subsection{What downstream papers may cite}

\begin{enumerate}
\item (Foundations, Hodge companion.) The $\Q$-linear projector
  $\Pi_\sigma \colon V_K \to V_K$ as in Definition~\ref{def:pisigma}:
  $\Q$-linear idempotent whose image is $\Fix(\sigma_V; V_K)$ and
  whose kernel is $V\delta$.
\item (Foundations, Hodge companion.) The fixed-part theorem
  (Theorem~\ref{thm:fixed-part}): the $\sigma_V$-fixed subspace of
  $V_K = V \otimes_\Q K$ is canonically $V$, for every
  finite-dimensional $\Q$-vector space $V$.
\item (Foundations, Hodge companion.) The basic properties of
  $\Pi_\sigma$ (Proposition~\ref{prop:pisigma-basic}): idempotence,
  image, kernel, coordinate formula.
\item (Foundations, Hodge companion.) The functoriality theorem
  (Theorem~\ref{thm:pisigma-functorial}): for every $\Q$-linear map
  $T \colon V \to W$, both $\sigma$ and $\Pi_\sigma$ are
  functorial under $T_K$.
\item (Foundations, Hodge companion.) The cohomological
  functoriality corollaries
  (Corollary~\ref{cor:cohomological} for pullback,
  Corollary~\ref{cor:cohomological-pushforward} for pushforward):
  cohomological maps defined over $\Q$ (in the precise sense of
  Definition~\ref{def:defined-over-Q}) commute with $\sigma_V$ and
  $\Pi_\sigma^V$ after scalar extension to $K$.
\end{enumerate}

\subsection{Recommended citation phrasing}

``By Theorem~3.2 and Corollaries~4.7 and~4.8 of
\cite{SigmaRationality}, the coefficientwise involution
$\sigma_{H^*(X, \Q)}$ on $H^*(X, \Q(\varphi)) \cong H^*(X, \Q)
\otimes_\Q \Q(\varphi)$ has fixed subspace $H^*(X, \Q) \otimes 1
\cong H^*(X, \Q)$, and the $\Q$-linear projector $\Pi_\sigma =
(\id + \sigma_{H^*(X, \Q)})/2$ commutes with the scalar extension
of every $\Q$-linear cohomological map: in particular, the singular
pullback $f^*$ of a continuous map (Corollary~\ref{cor:cohomological}),
and any externally supplied $\Q$-linear cohomological pushforward
$f_*$ (Corollary~\ref{cor:cohomological-pushforward}, whose
existence is a hypothesis of that corollary, not a conclusion).''

\subsection{What this paper does not provide}

The following are \emph{not} proved here and must be justified
separately in any downstream paper that needs them:

\begin{enumerate}
\item Any step from ``$\sigma$-fixed'' to ``algebraic'' (i.e.\ any
  Hodge-Conjecture-type step).
\item Any cascade-embeddability or universality statement for
  varieties.
\item Any non-coefficientwise $\sigma$-action on configurations,
  dynamics, metrics, gauge fields, cohomology of non-$\Q$-defined
  sheaves, or motivic structures.
\item Any statement over $\Z$, over torsion coefficients, or about
  integral lattice structures beyond the generic scalar-extension
  structure here.
\item Any relation between $\Pi_\sigma$ and Langlands, periods,
  motives, or Hodge-theoretic non-abelian duality.
\end{enumerate}

\subsection{Relation to the Hodge Millennium problem}

The ``$\sigma$-fixed = rational'' step that appears informally in
earlier cascade Hodge drafts is exactly
Theorem~\ref{thm:fixed-part} of this paper, but the ``rational
$\implies$ algebraic'' step (i.e.\ the Hodge Conjecture proper) is
entirely orthogonal to the content here. A reader who sees $\Pi_\sigma$ cited in a downstream paper
should not interpret that citation as evidence for the Hodge
Conjecture; it is only evidence that the specific algebraic
bookkeeping step invoked there is rigorous.

%=====================================================================
\appendix
\section{Notation Audit}
\label{app:notation}
%=====================================================================

\begin{center}
\begin{tabular}{lll}
\hline
Symbol & Meaning & Formula / reference \\
\hline
$\varphi$ & golden ratio & $\varphi = (1+\sqrt{5})/2$ \\
$\bar\varphi$ & Galois conjugate of $\varphi$ & $\bar\varphi = 1 - \varphi = -1/\varphi$ \\
$\delta$ & $\sigma$-odd generator & $\delta = 2\varphi - 1 = \sqrt{5}$ \\
$K$ & the field $\Q(\varphi)$ & $K = \Q \oplus \Q\delta$ (Lemma~\ref{lem:K-decomp}) \\
$\sigma$ & Galois involution of $K/\Q$ & $\sigma(a + b\varphi) = a + b\bar\varphi$; $\sigma(\delta) = -\delta$ \\
$V_K$ & scalar extension of $V$ to $K$ & $V_K = V \otimes_\Q K$ (Def.~\ref{def:VK}) \\
$V\delta$ & $\sigma$-anti-fixed summand & $V\delta = V \otimes_\Q \Q\delta \subset V_K$ \\
$\iota_V$ & inclusion $V \hookrightarrow V_K$ & $\iota_V(v) = v \otimes 1$ (Def.~\ref{def:inclusion}) \\
$\sigma_V$ & coefficientwise $\sigma$ on $V_K$ & $\sigma_V = \id_V \otimes \sigma$ (Def.~\ref{def:sigmaV}) \\
$\Pi_\sigma$ & fixed-part projector & $\Pi_\sigma = (\id + \sigma_V)/2$ (Def.~\ref{def:pisigma}) \\
$\Fix$ & fixed subspace & $\Fix(\sigma_V; V_K) = \{\omega : \sigma_V(\omega) = \omega\}$ \\
$T_K$ & scalar extension of $T$ & $T_K = T \otimes \id_K$ (Def.~\ref{def:TK}) \\
\hline
\end{tabular}
\end{center}

\bigskip

\noindent\textbf{Coordinate formulas.} For $\omega \in V_K$ with
$\delta$-basis decomposition $\omega = v_0 + v_1 \delta$ ($v_0, v_1
\in V$):
\[
  \sigma_V(\omega) = v_0 - v_1 \delta, \qquad
  \Pi_\sigma(\omega) = v_0.
\]
Equivalently, in the $\varphi$-basis $\omega = u_0 + u_1 \varphi$
($u_0, u_1 \in V$), with $v_0 = u_0 + u_1/2$ and $v_1 = u_1/2$:
\[
  \sigma_V(\omega) = (u_0 + u_1) - u_1 \varphi, \qquad
  \Pi_\sigma(\omega) = u_0 + u_1/2.
\]

\begin{thebibliography}{99}

\bibitem{LangAlgebra}
S.~Lang, \emph{Algebra}, 3rd ed., Graduate Texts in Mathematics
\textbf{211}, Springer, 2002. Galois theory: Ch.~V;
semilinear algebra: Ch.~XVI \S 7.

\bibitem{AtiyahMacdonald}
M.~F.~Atiyah and I.~G.~Macdonald, \emph{Introduction to Commutative
Algebra}, Addison-Wesley, 1969. Tensor products and scalar
extension: Ch.~2.

\bibitem{Hatcher}
A.~Hatcher, \emph{Algebraic Topology}, Cambridge University Press,
2002. $H^*(\CP^n, \Q)$: Theorem 3.19.

\bibitem{GriffithsHarris}
P.~Griffiths and J.~Harris, \emph{Principles of Algebraic Geometry},
Wiley, 1978. Cohomology of abelian varieties: Ch.~2.

\bibitem{SigmaRationality}
(This paper.)
\emph{Cascade--Classical Correspondence I: Sigma-Fixed Rationality
over $\Q(\varphi)$}.

\end{thebibliography}

\end{document}
